\scanpage{104}

\mspage{110}

\margin{$\unsure{\ast|\times}$}

La première assertion doit résulter des formules

\[
\begin{aligned}
\deleted{\unsure{\mathrm{Hom}}}\qquad H^{\unsure{\cdot|i}}_{\mathrm{HOT}}(Z,a_Z)&\simeq\underline{\mathcal H\mathrm{om}}_{\mathrm{Hot}}(\mathrm{hot}(Z),a)\\
(\text{et aussi }\quad H_{\unsure{\cdot|i}}^{\mathrm{HOT}}(Z,\unsure{a|e}_{\unsure{Z|z}})&\simeq\mathrm{hot}(Z)\times a\quad ,
\end{aligned}
\]

\deleted{\unsure{prenant} le $\pi_0$}

donc : si $(**)\ \mathrm{hot}(X'')\xrightarrow{\sim}\mathrm{hot}(X')$ est iso, alors $H^{\cdot}_{\mathrm{Hot}}(X^{\unsure{'|''}},a_{X'})\longrightarrow H^{\cdot}_{\mathrm{HOT}}(X'',a_{X''})$ est iso, puisque

c'est déduit de $(*)$ en prenant $\underline{\mathcal H\mathrm{om}}_{\mathrm{Hot}}(\cdot,a)$. Inversement, si $(*)$ est tjrs iso, \deleted{par exp} \margin{$\longrightarrow$}exprimant les deux \unsure{mb's|mbres} comme des $\underline{\mathcal H\mathrm{om}}_{\mathrm{Hot}}$ et prenant les $\pi_0$, on trouve que

\[
\mathrm{Hom}_{\mathrm{Hot}}(\mathrm{hot}(X'),a)\xrightarrow{\sim}
\mathrm{Hom}_{\mathrm{Hot}}(\mathrm{hot}(X''),\unsure{\cdot|a})
\]

est iso pour tout $a$, donc $(**)$ est iso.

Mais le fait que ça implique que $f''$ soit une \inserted{HOT-}équiv. \deleted{\unsure{$\mathbb D$}} coh. complète, i.e. que la relation $(*)$ reste vraie pour un $\unsure{\xi|\zeta}\in\mathbb D(X')$ \uline{loc. \unsure{constant}} sur $X'$, me paraît une \deleted{chose} propriété délicate du dérivateur HOT, sur laquelle il me faudra revenir.

Donc le lemme est loin d'être établi, faute d'avoir seulement construit le dérivateur HOT, sans $=$ parler d'en établir des propriétés spéciales délicates ! Mais je n'ai aucune raison de douter du lemme, il fait partie \inserted{de la kyrielle} des intuitions de base que je \deleted{\unsure{avais}} dois \deleted{\unsure{à}} fonder solidement.

% archive sheet number: 104
